\subsection{规律性生活的健康获益：证据与流言}

上一小节提到"规律的性生活与较低的死亡风险相关"。这里把"性对身体的好处"和网上流传的种种说法分开讲——\textbf{获益要讲证据，流言要划清界线}。

\vspace{0.5em}
\noindent\textbf{有证据支持的获益}：

\begin{itemize}
	\item \textbf{心血管}：一次典型的性生活相当于快走或爬楼梯的强度（代谢当量约 3--5 METs），规律进行对心肺功能是温和的锻炼；大规模队列研究显示，频率较高的男性\textbf{心血管事件风险更低}（注意：是"健康者性更多"与"性更多者更健康"的双向关系，不宜简单因果化）
	\item \textbf{免疫}：有研究（Pennsylvania 的经典研究）发现规律性生活的大学生唾液免疫球蛋白 A（IgA）水平更高；此外性高潮后催乳素升高与睡眠改善也与免疫调节间接相关
	\item \textbf{睡眠与情绪}：性高潮释放的催乳素、催产素与内啡肽有助于入睡和缓解焦虑——这也是"睡前性生活"被反复验证的实用价值
	\item \textbf{盆底与产后恢复}：规律、无痛的性生活提示盆底功能状态良好；性活动中的盆底收缩本身就是一种"用进废退"
	\item \textbf{关系与寿命}：亲密关系质量是全因死亡率的重要预测因子；性是维系成人依恋的重要通道——其获益\textbf{主要通过"关系质量"间接兑现}
\end{itemize}

\vspace{0.5em}
\noindent\textbf{常见流言与夸大说法，逐条澄清}：

\begin{itemize}
	\item \textbf{"做一次燃烧 300 大卡，性爱是最好的减肥运动"}——言过其实。实测平均一次约消耗 \textbf{85--150 大卡}（3--7 分钟的中低强度活动），"躺赢减肥"不存在；说它"运动强度接近爬两层楼"更准确
	\item \textbf{"性生活多，皮肤就好、还能防脱发"}——没有可靠证据。性高潮后的短暂"红润"来自血管扩张，不会改写皮肤老化或毛囊命运
	\item \textbf{"精液是营养品，吞了美容"}——精浆 95\% 以上是水，蛋白质与微量元素含量微不足道（详见"口交中的射精、吞咽与精液安全"一节）
	\item \textbf{"射精次数固定，年轻时用多了老了就不行"}——"一生射精配额"没有任何生理学依据；规律射精（每月 21 次以上的流行病学阈值）反而与\textbf{较低的前列腺癌风险}相关（大队列研究提示，机制仍在研究）
	\item \textbf{"性生活能治百病"}——性是健康的"信号灯"与"助推器"，不是治疗手段；把性当作药，反而会推迟就医
\end{itemize}

\vspace{0.5em}
\noindent\textbf{一个重要的反向提醒}：性功能的突然改变（勃起变差、性欲骤降、高潮消失）常是\textbf{全身疾病的早期信号}（血管、内分泌、抑郁、药物）——与其纠结"做多少次最养生"，不如把"性"当作每年一次的免费自检：\textbf{一直正常，说明底子不错；突然变差，该去体检了}。

\begin{tcolorbox}[colback=pink!5!white,colframe=red!50!black,title=贴心小叮咛]
性生活对健康的贡献，\textbf{七分来自关系与心情，三分来自身体活动}。它值得被认真对待，但不该被神化——规律、愉悦、双方都舒服的性，本身就是目的，顺便的"健康红利"只是额外的馈赠。
\end{tcolorbox}

